The hardware-autonomous I2C target (\peripheral{I2CTx}) is an I2C slave engine that handles the wire protocol entirely in hardware: it matches its 7-bit address (with a per-bit wildcard mask and an optional general-call response), receives and transmits one byte at a time with ready/empty status, stretches the clock for lossless flow control, and flags every framing event (START, STOP, repeated-START, NACK) and a stuck-bus condition. It complements the software-serviced slave-mode registers of \peripheral{I2C0}/\peripheral{I2C1}: where those cores expose the slave state to firmware for per-byte servicing, \peripheral{I2CT0} needs no bit-banging --- it answers its address, ACKs, buffers a byte, and raises an interrupt, all in hardware. The whole engine runs on the free-running fabric clock (MCLK); the SDA and SCL pins are each sampled through an independent two-flop synchronizer, so no flip-flop is ever clocked by a bus pin (the deliberate contrast with the pad-clocked slave path of \peripheral{I2C0}). The main features are listed below:

\begin{itemize}
	\item 7-bit address match with a wildcard mask (\register{I2CTSAD}/\register{I2CTSADM}, a mask bit makes that address bit a don't-care) and an optional general-call response (\register{I2CTGCEN})
	\item Byte-at-a-time receive (\register{I2CTxRX} + \register{I2CTRXF}) and transmit (\register{I2CTxTX} + \register{I2CTTXE}), MSB-first per the I2C convention, with a side-effect-free received-byte read
	\item Hardware clock stretching (\register{I2CTCSEN}): the target holds SCL low after an ACK until firmware services the transfer, giving lossless flow control; with stretching disabled the target never holds SCL and firmware must service within one bit-time
	\item Full framing detection --- \register{I2CTSTOPF}, \register{I2CTRSTARTF}, \register{I2CTNACKF}, \register{I2CTAMF}/\register{I2CTGCF} --- plus a read-only bus-busy (\register{I2CTBUSY}) and transfer-direction (\register{I2CTTM}) status
	\item Receive-overrun protection: a byte arriving while the previous one is unread sets \register{I2CTOVF}, is auto-NACKed and dropped
	\item A configurable stuck-SCL watchdog (\register{I2CTxWDG}): a timeout on SCL held low while busy sets \register{I2CTERRF} and aborts the transaction; the reset value 0 disables it (default-off)
	\item Sticky write-1-to-clear status flags and two combined interrupts (vector 122 address/error, vector 123 tx-ready/rx-full)
\end{itemize}

\subsection{Address Match, Direction, and Framing}

On a START the target shifts in the address byte MSB-first and matches it against \register{I2CTSAD} masked by \register{I2CTSADM} (a mask bit is a don't-care); with \register{I2CTGCEN} set it also answers the general-call address (0x00, write direction) and sets \register{I2CTGCF} alongside \register{I2CTAMF}. A match drives the address ACK (SDA pulled low --- the target never drives SDA high, open-drain) and latches the transfer direction into \register{I2CTTM} (1 = target-transmitter / host read, 0 = target-receiver / host write); a mismatch is silently ignored. \register{I2CTBUSY} tracks the transaction from START to STOP and distinguishes an initial START from a repeated-START (a START while busy sets \register{I2CTRSTARTF} and re-enters the address phase without dropping busy --- the classic write-address-then-read reversal). Every framing event releases the bus immediately, so the target never holds SDA or SCL across a boundary.


The wire protocol itself --- START, address and direction, ACK, data, STOP --- is the same one shown in Figure \ref{fig:i2c-transaction} in the \peripheral{I2Cx} chapter.

\subsection{Receive, Transmit, and Clock Stretching}

On a host write the target shifts the byte in, and if the receive buffer is free it ACKs, latches the byte into \register{I2CTxRX} and sets \register{I2CTRXF}; firmware learns via the vector-123 interrupt, reads \register{I2CTxRX} (side-effect-free) and writes 1 to \register{I2CTRXF} to free the buffer. On a host read the target needs bytes from firmware: \register{I2CTTXE} asserts when the transmit buffer is empty, firmware loads \register{I2CTxTX} (a write clears \register{I2CTTXE} in the same cycle it loads a byte, so a poll right after the write is never blind), and the byte shifts out MSB-first. When clock stretching is enabled (\register{I2CTCSEN}) the target holds SCL low after the ACK until firmware services the transfer --- the host, waiting to raise SCL for the next bit, simply stalls; release is synchronous in MCLK, which is safe because the host is already stopped. With stretching disabled the target never holds SCL: firmware must keep up within one bit-time or accept an overrun (\register{I2CTOVF}, auto-NACK) on receive, or a defined stale 0xFF on transmit.

\subsection{Bus Speed and the MCLK Ratio}

The block is purely edge-reactive --- it never counts raw clock cycles for protocol timing --- so nothing in it scales with the actual bus frequency; only the synchronizer and edge-detect latency (a few MCLK cycles per SDA/SCL event) sets a floor on how fast SCL may run. The guaranteed floor is $f_\mathrm{SCL} \le \mathrm{MCLK}/24$: at the nominal 24\,MHz MCLK both Standard-mode (100\,kHz) and Fast-mode (400\,kHz) are guaranteed with ample margin. Because the engine clocks off the free-running fabric clock, its timing is immune to clock reconfiguration, and software must \emph{not} apply the ``write \register{\register{SYSCLKCR}} to 0 first'' rule that the SMCLK-domain \peripheral{I2C0}/\peripheral{I2C1} cores need --- \peripheral{I2CT0} is unaffected by SMCLK reconfiguration.

\subsection{Interrupts and Shared Access on \AsicNameForUserGuide}

The two combined interrupts are $(\register{I2CTAMF} \vee \register{I2CTGCF} \vee \register{I2CTOVF} \vee \register{I2CTNACKF} \vee \register{I2CTSTOPF} \vee \register{I2CTRSTARTF} \vee \register{I2CTERRF}) \wedge \register{I2CTAEIE}$ on vector 122 (address/error) and $(\register{I2CTRXF} \vee \register{I2CTTXE}) \wedge \register{I2CTDATAIE}$ on vector 123 (tx-ready/rx-full), both combinational and never latched; the service routine reads \register{I2CTxSR}, demultiplexes the source and clears whichever event flags fired before returning. In configurations that include it, \peripheral{I2CT0} lives in the shared peripheral window behind the multi-core arbiter (see Section \ref{s:multicore}) in the mutex-bank page, sub-slot 10 at \texttt{0x6A00}, so any hart can use it. Its two interrupt sources occupy vectors 122 and 123, above the external-interrupt (meip) delivery vector fixed at 85 and above every other shared-peripheral source.

The block adds no pins: it shares the open-drain \texttt{SDA0}/\texttt{SCL0} pads with \peripheral{I2C0} through a wired-AND merge on the drive-low (DIR) plane --- the SCL merge realizes the target's clock stretching and the SDA merge its ACK/data driving --- so the two engines can share one bus, including the on-die host-side loopback used for self-test. When \peripheral{I2CT0} is not instantiated, those pads carry only \peripheral{I2C0} exactly as before.
